\section{From multivariate temporal search to a conditional effect query}
\label{s5:discovery}
For deposit pricing, balances, lending and market conditions, a useful
discovery workflow must allow lagged and contemporaneous relations. PCMCI+
uses time order and parent selection to address autocorrelation under causal
sufficiency; its contemporaneous output can retain unoriented edges
\citep{runge2020}. LPCMCI allows hidden common causes within its stated
time-series graphical model. A partial-correlation test supplies a linear
Gaussian conditional-independence test, not generic nonlinear independence.
Stationarity, causal Markov, faithfulness, a sufficient lag window and test
validity remain separate premises. Choosing a flexible test changes sample
requirements and computation as well as the null it tests.

The graphical target is itself a method choice. The temporal ancestral-graph
characterization distinguishes the latent projection on a finite window from
its stationarified representation, with potentially different orientation
information \citep[Sections 4.5--4.6 and Propositions 5.2--5.7]{gerhardus2021}.
Its conceptual complete enumeration may be computationally infeasible, and
the inspected 2021 preprint explains why then-current LPCMCI does not recover
every invariant of the stronger time-series graph class. This is not a
claim about every later version. TS-ICD supplies a distinct ordering of
temporal and contemporaneous search under its stationary SVAR/oracle
assumptions \citep{rohekar2023}. Faster orientation under one graph class
is not evidence of completeness under another. The preserved author versions
and later journal identities are recorded separately.

The selected runtime comparison uses Tigramite 5.2.10.1, PCMCI+ and LPCMCI,
with $\alpha=.01$, six variables, 1,600 dates and a maximum lag of two
\citep{runtime_tigramite_scope}. LPCMCI's conditioning search is capped at
three; this is an explicit computational restriction, so no unrestricted
oracle-completeness claim is made. The generator has stable lagged effects,
two contemporaneous triangular relations and independent Gaussian
innovations. A challenger adds an unobserved autocorrelated common cause
of two series. Another searches only lag one while the true mechanism
retains lag two. All graph endpoint strings, conditional-independence
statistics, $p$-matrices and timings are preserved over 20 repetitions.

\begin{center}\small
\begin{tabular}{lrrrr}\toprule
&\multicolumn{2}{c}{PCMCI+}&\multicolumn{2}{c}{LPCMCI}\\
Regime&Precision&Recall&Precision&Recall\\\midrule
Sufficient, lag two&.993&1.000&.971&1.000\\
Hidden common cause&.865&1.000&.815&1.000\\
Omitted second lag&.980&.929&.917&.929\\\bottomrule
\end{tabular}\end{center}
These are descriptive adjacency comparisons with the observed structural
DAG. In the hidden case that DAG adjacency is not the PAG target, so the
numbers must not be read as PAG accuracy or a ranking of latent discovery
correctness. Recall counts true edges outside the requested lag window;
in-window recall alone is one and conceals the omitted-lag defect. Finite
$p$-values and edge-selection frequencies are not graph-confidence levels.

\subsection{A supported transformation and its preserved uncertainty}
The local effect adapter accepts a supplied sufficient temporal graph with
directed lagged edges and directed or unoriented contemporaneous edges. For
at most seven variables it enumerates contemporaneous topological orders
that preserve directed edges and introduce no unrepresented unshielded
collider, including colliders involving lagged parents. It refuses conflicting,
bidirected and partially ancestral endpoints. This is a bounded sufficient
graph operation, not a general PAG completion algorithm.

For each permitted order fit the declared linear structural equations on a
fresh independent 3,000-date sample:
\[
 (I-B_0)X_t=\sum_{\ell=1}^p B_\ell X_{t-\ell}+\varepsilon_t.
\]
Only parents present in that candidate graph enter each equation. The
companion matrix built from $(I-B_0)^{-1}B_\ell$ propagates an additive
equation impulse $(I-B_0)^{-1}e_a$. Query outcome $j$ at horizon $h$ gives
the corresponding coordinate of the propagated vector. The result contains
each order, stability diagnostic and curve, plus a pointwise envelope.
The tested PCMCI+ sufficient graph admits 180 compatible orders, many
differing only in the order of unrelated variables; these are not 180
distinct causal models. LPCMCI output retaining richer marks is refused
by this adapter. The hidden-cause case is outside its sufficiency model,
not proof that every hidden-variable effect is unidentified.

The output is conditional on the selected graph and linear structural
restrictions. A fresh estimation sample prevents reusing search noise in the
coefficient fit, but does not guarantee the selected graph contains the
truth. Consequently the displayed envelope is structural-order sensitivity,
not a confidence band. The earlier finite covariance-region construction
in Section~\ref{s4:graph-confidence} provides a distinct confidence route
under its own bounded model; its guarantee cannot be attached to this search.

A two-series example makes the remaining decision ambiguity material.
Let reduced-form innovations have covariance
$\left(\begin{smallmatrix}1&.6\\.6&1.36\end{smallmatrix}\right)$ and
autoregression $.5I$. One recursive model has a contemporaneous
price-to-balance coefficient .6; a reverse recursive model fits the same
Gaussian innovation law and gives zero balance response to the price
equation impulse. With 8,000 simulated observations the supported order
envelope at impact is $[0,.6096]$, with approximately geometric propagation.
Neither more observations of that unchanged Gaussian law nor a lower search
threshold chooses the causal direction. A documented timing restriction or
an externally randomized price perturbation would distinguish these models.
If the action decision needs a response above .3, the existing evidence does
not make it robust across the two models. This is an actionable request for
additional information, not a decorative ambiguity label.

The practical hierarchy is therefore PCMCI+ for a defensible sufficient
stationary temporal model, LPCMCI as a hidden-confounding challenger with
retained endpoint semantics, and structural identification/effect sets only
where the query's premises justify them. Increasing lag windows can add
identifying information at the population level but increases finite-sample
conditioning cost. Interventions, reliable prior ordering and additional
measurements should be chosen for the effect ambiguity they resolve. A
large discovery runtime cannot by itself convert a search graph into a
bank policy recommendation.
